Un nombre $n \geqslant 2$ de segments se trouvent dans le plan de sorte que deux segments quelconques se
coupent en un point différent de leurs extrémités, mais sans que trois d'entre eux possèdent un point commun.
Pour chaque segment, Geoff doit choisir une extrémité et y placer une grenouille tournée vers l'autre
extrémité. Ensuite, il doit claquer des mains, au total $n - 1$ fois. Chaque fois qu'il claque des mains,
chaque grenouille saute immédiatement en avant sur son segment jusqu'au point d'intersection suivant. Les
grenouilles ne changent jamais de direction.

Geoff souhaite placer les grenouilles de sorte qu'il n'y ait jamais plus d'une grenouille à chaque point
d'intersection.
\begin{enumerate}
  \item[(a)] Montrer que si $n$ est impair, Geoff peut toujours trouver un placement des grenouilles qui
  satisfasse son souhait.
  \item[(b)] Montrer que si $n$ est pair, aucun placement des grenouilles ne satisfait son souhait.
\end{enumerate}
